% Source: https://artofproblemsolving.com/community/c3h2482455_fun_problems_
% Stable identity: aops-topic-2482455-q01
Take $I=(0,0),V=(2,0),Y=(1,\sqrt3),A=(1,0)$. The ray $AS$ has equation $y=1-x$, while $IY$ has equation $y=\sqrt3x$. Thus $S$ has $x$-coordinate $\dfrac{\sqrt3-1}{2}$. Consequently $\dfrac{YS}{IY}=\dfrac{3-\sqrt3}{2}$. The altitude from $A$ to line $IY$ is half that from $V$, so the area ratio is $\dfrac{3-\sqrt3}{4}$. Hence $\boxed{a+b+c=3+3+4=10}$.
\begin{asy}
usepackage("amsmath");
size(7cm);
pair I=(0,0),V=(2,0),Y=(1,sqrt(3)),A=(1,0),PtS=((sqrt(3)-1)/2,(3-sqrt(3))/2);draw(I--V--Y--cycle);draw(A--PtS--Y--A,blue);label("$I$",I,SW);label("$V$",V,SE);label("$Y$",Y,N);label("$A$",A,S);label("$S$",PtS,W);
\end{asy}
